% Part: proof-theory
% Chapter: sequent-calculus
% Section: invertibility

\documentclass[../../../include/open-logic-section]{subfiles}

\begin{document}

\olfileid{pt}{seq}{inv}

\olsection{Invertibility of Rules}

\begin{defn}
A rule~$R$,
\[
\AxiomC{$S_1$}
\AxiomC{$\dots$}
\AxiomC{$S_n$}
\RightLabel{$R$}
\TrinaryInfC{$S$}
\DisplayProof
\]
is \emph{invertible} in a system if, whenever $\Proves S$ then
$\Proves[h_i] S_i$ for $i = 1$, \dots,~$n$. It is
\emph{!!{height}-preserving invertible} (or \emph{!!{hp}-invertible})
in a system if, whenever $\Proves_n S$ then $\Proves[n] S_i$ for $i =
1$, \dots,~$n$.
\end{defn}

\begin{lem}\ollabel{lem:G3c-invert}
The rules of~\Log{G3c} for $\lnot$, $\land$, $\lor$, and $\lif$ are
!!{height}-preserving invertible, that is:
\begin{enumerate}
  \item \RightR{\lnot}: If $\Log{G3c} \Proves[n] \Gamma \Sequent
  \Delta, \lnot !A$, then $\Log{G3c} \Proves[n] !A, \Gamma \Sequent
  \Delta$.
  \item \LeftR{\lnot}: If $\Log{G3c} \Proves[n] \Gamma, \lnot !A
  \Sequent \Delta$, then $\Log{G3c} \Proves[n] \Gamma \Sequent \Delta,
  !A$.
  \item \RightR{\land}: If $\Log{G3c} \Proves[n] \Gamma \Sequent
  \Delta, !A \land !B$, then $\Log{G3c} \Proves[n] \Gamma \Sequent
  \Delta, !A$ and $\Log{G3c} \Proves[n] \Gamma \Sequent
  \Delta, !B$.
  \item \LeftR{\land}: If $\Log{G3c} \Proves[n] \Gamma, !A \land !B
  \Sequent \Delta$, then $\Log{G3c} \Proves[n] !A, !B, \Gamma \Sequent
  \Delta$.
  \item \RightR{\lor}: If $\Log{G3c} \Proves[n] \Gamma \Sequent
  \Delta, !A \lor !B$, then $\Log{G3c} \Proves[n] \Gamma \Sequent
  \Delta, !A, !B$.
  \item \LeftR{\lor}: If $\Log{G3c} \Proves[n] !A \lor !B, \Gamma \Sequent
  \Delta$, then $\Log{G3c} \Proves[n] !A, \Gamma \Sequent
  \Delta$ and $\Log{G3c} \Proves[n] !B, \Gamma \Sequent
  \Delta$.
  \item \RightR{\lif}: If $\Log{G3c} \Proves[n] \Gamma \Sequent
  \Delta, !A \lif !B$, then $\Log{G3c} \Proves[n] !A, \Gamma \Sequent
  \Delta, !B$.
  \item \LeftR{\lif}: If $\Log{G3c} \Proves[n] !A \lif !B, \Gamma
  \Sequent \Delta$, then $\Log{G3c} \Proves[n] \Gamma \Sequent \Delta,
  !A$ and $\Log{G3c} \Proves[n] !B, \Gamma \Sequent \Delta$.
\end{enumerate}
\end{lem}

\begin{proof}
We have to show that, for any rule~$R$ of~\Log{G3c}, if there is
!!a{proof}~$\pi$ of the conclusion~$S$ of~$R$, then there are also
!!{proof}s~$\pi_i$ of the premises $S_i$ of~$R$, with $\pheight{\pi_i}
\le \pheight{\pi}$. Consider the rule \LeftR{\land}: suppose there is
!!a{proof}~$\pi$ with end-sequent of the form $!A \land !B, \Gamma
\Sequent \Delta$. We have to show that there is !!a{proof} $\pi'$ of
$!A, !B, \Gamma \Sequent \Delta$ with $\pheight{\pi'} \le
\pheight{\pi}$. We do this by induction on $n = \pheight{\pi}$.

If $n = 0$, $\pi$ consists of just an axiom: $!A \land !B, \Gamma
\Sequent \Delta$. In \Log{G3c}, axioms are of the form $!C, \Gamma
\Sequent \Delta, !C$ with $!C$ atomic. In other words, $!C$ can't be
$!A \land !B$. If $!A \land !B, \Gamma \Sequent \Delta$ is an axiom,
then some atomic $!C$ must be !!a{element} of both $\Gamma$
and~$\Delta$. But then $!A, !B, \Gamma \Sequent \Delta$ is also an
axiom, and we have found our !!{proof}~$\pi'$ (which is also of
!!{height}~$0$).

If $n>0$ it ends in an inference. We have to prove the claim for~$n$,
but may assume that it is true for any !!{proof} of height~$<n$, even
if the context is different. In other words, for any $k<n$ and any
$\Pi$ and $\Lambda$, if $!A \land !B, \Pi \Sequent \Lambda$ has
!!a{proof} of !!{height}~$k$, then there is !!a{proof} of
!!{height}~$\le k$ of $!A, !B, \Pi \Sequent \Lambda$.

We have two cases: either $!A \land
!B$ on the left is principal, or it is not. If it is principal, the
last inference is \LeftR{\land} and the immediate
sub-!!{proof}~$\pi_1$ ends in $!A, !B, \Gamma \Sequent \Delta$. The
result holds in this case as $\pheight{\pi_1} < \pheight{\pi}$.

If $!A \land !B$ is not principal, some other !!{formula} is and $!A
\land !B$ belongs to the context of the last inference. Then the
inductive hypothesis applies to the immediate sub-!!{proof}s of the
last rule~$R$, which yields !!a{proof}~$\pi_1'$ or two
!!{proof}s~$\pi_1'$ and~$\pi_2'$ of !!{height} no greater than the
height(s) of the immediate sub-!!{proof}s of~$\pi$. The end-sequents
of $\pi_1'$ and $\pi_2'$ contain $!A, !B$ in the context on the left
instead of $!A \land !B$. Apply the rule~$R$ to obtain the
proof~$\pi'$. For instance, suppose $\pi$ ends in~\LeftR{\lif}:
\[
\AxiomC{}
\RightLabel{$\pi_1$}
\Deduce$!A \land !B, \Gamma' \fCenter \Delta, !C$
\AxiomC{}
\RightLabel{$\pi_2$}
\Deduce$!A \land !B, \Gamma', !D \fCenter \Delta$
\RightLabel{\LeftR{\lif}}
\BinaryInf$!A \land !B, \Gamma', !C \lif !D \fCenter \Delta$
\DisplayProof
\]
Here $!C \lif !D$ on the left is principal, and the context on the
left is $!A \land !B, \Gamma'$ (i.e., $\Gamma = \Gamma', !C \lif !D$).
The inductive hypothesis yields !!{proof}s $\pi_1'$ and $\pi_2'$ of 
$!A \land !B, \Gamma' \Sequent \Delta, !C$ and $!A \land !B, \Gamma',
!D \Sequent \Delta$, respectively. These two sequents can again
serve as premises for the \LeftR{\lif} rule to get~$\pi'$:
\[
\AxiomC{}
\RightLabel{$\pi_1'$}
\Deduce$!A, !B, \Gamma' \fCenter \Delta, !C$
\AxiomC{}
\RightLabel{$\pi_2'$}
\Deduce$!A, !B, \Gamma', !D \fCenter \Delta$
\RightLabel{\LeftR{\lif}}
\BinaryInf$!A, !B, \Gamma', !C \lif !D \fCenter \Delta$
\DisplayProof
\]
We have 
\begin{align*}
\pheight{\pi} & = \max(\pheight{\pi_1}, \pheight{\pi_2}) + 1\\
\pheight{\pi'} & = \max(\pheight{\pi_1'}, \pheight{\pi_2'}) + 1
\intertext{by definition, and}
\pheight{\pi_1'} & \le \pheight{\pi_1}\\
\pheight{\pi_2'} & \le \pheight{\pi_2}
\intertext{by inductive hypothesis. So,}
\pheight{\pi'} & \le \pheight{\pi}.
\end{align*}
\end{proof}

\begin{prob}
Give the proof of \olref{lem:G3c-invert} for \RightR{\lif}.
\end{prob}

\begin{lem}\ollabel{lem:invert-quant}
The rules \RightR{\lforall} and \LeftR{\lexists} are
!!{height}-preserving invertible in the following sense:
\begin{enumerate}
  \item If $\Log{G3c} \Proves[n] \Gamma \Sequent \Delta,
  \lforall[x][!A(x)]$ then $\Log{G3c} \Proves[n] \Gamma \Sequent
  \Delta, !A(c)$, and
  \item if $\Log{G3c} \Proves[n] \lexists[x][!A(x)], \Gamma \Sequent
  \Delta$ then $\Log{G3c} \Proves[n] !A(c), \Gamma \Sequent \Delta$,
\end{enumerate}
for any $c$ not occuring in $\Gamma$, $\Delta$, and $\lforall[x][!A(x)]$
or $\lexists[x][!A(x)]$, respectively.
\end{lem}

\begin{proof}
  We show (1) and leave (2) as an exercise. The argument is similar to
  the proof of \olref{lem:G3c-invert}. In the inductive step, there
  are again two cases. The case where $\lforall[x][!A(x)]$ is
  principal again is quick: the premise of the last
  $\RightR{\lforall}$ inference is of the required form $\Gamma
  \Sequent \Delta, !A(a)$ and the immediate sub-!!{proof}~$\pi_1$
  ending in it has $\pheight{\pi_1} < \pheight{\pi}$. Moreover, the
  eigenvariable condition gives us that $a$ does not occur in
  $\Gamma$, $\Delta$, or $\lforall[x][!A(x)]$. Since we may assume that
  $\pi_1$ is regular, $\Subst{\pi_1}{c}{a}$ is !!a{proof} of the same
  !!{height} of $\Gamma \Sequent \Delta, !A(c)$ by
  \olref[qua]{lem:subst-regular}.

  In the second case, $\lforall[x][!A(x)]$ is not principal in the
  last inference; some other formula is. We again consider a case by
  way of example. Suppose the last inference is~\RightR{\land}. Some
  !!{formula} in $\Delta$ is of the form~$!B \land !C$ and is principal.
  The !!{proof}~$\pi$ ends in
  \[
\AxiomC{}
\RightLabel{$\pi_1$}
\Deduce$\Gamma \fCenter \Delta', \lforall[x][!A(x)], !B$
\AxiomC{}
\RightLabel{$\pi_2$}
\Deduce$\Gamma \fCenter \Delta', \lforall[x][!A(x)], !C$
\RightLabel{\RightR{\land}}
\BinaryInf$\Gamma \fCenter \Delta', \lforall[x][!A(x)], !B \land !C$
\DisplayProof
\]
where $\Delta = \Delta', !B \land !C$. Let $c$ be !!a{constant} that
does not appear in $\Delta$; then clearly it also does not occur in
$\Delta', !B$ or in $\Delta', !C$. So the inductive hypothesis applies
to $\pi_1$ and~$\pi_2$; we have !!{proof}s~$\pi_1'$ and $\pi_2'$ with
$\pheight{\pi_1'} \le \pheight{\pi_1}$ and $\pheight{\pi_2'} \le
\pheight{\pi_2}$. The required !!{proof}~$\pi'$ of $\Gamma \Sequent
\Delta, !A(c)$ is:
\[
\AxiomC{}
\RightLabel{$\pi_1'$}
\Deduce$\Gamma \fCenter \Delta', !A(c), !B$
\AxiomC{}
\RightLabel{$\pi_2'$}
\Deduce$\Gamma \fCenter \Delta', !A(c), !C$
\RightLabel{\RightR{\land}}
\BinaryInf$\Gamma \fCenter \Delta', !A(c), !B \land !C$
\DisplayProof
\]
with $\pheight{\pi'} = \max(\pheight{\pi_1'}, \pheight{\pi_2'})+1 \le \pheight{\pi}$.
\end{proof}

\Olref{lem:G3c-invert} plays an important role in the proof of the
cut-elimination theorem. But it can also be used to show the
following:

\begin{prop}\ollabel{prop:G3c-cont-adm}
The contraction rules \LeftR{\Contraction} and \RightR{\Contraction}
of~\Log{G1c} are !!{height}-preserving admissible in~\Log{G3c}.
\end{prop}

\begin{proof}
We give the argument for \RightR{\Contraction} and
\LeftR{\Contraction} at the same time. Suppose $\pi$ is !!a{proof} of
$\Gamma \Sequent \Delta, !A, !A$ or of $!A, !A, \Gamma \Sequent
\Delta$ in~\Log{G3c}. We have to show that there is also a
\Log{G3c}-!!{proof}~$\pi'$ of $\Gamma \Sequent \Delta, !A$ or $!A,
\Gamma \Sequent \Delta$, respectively, with $\pheight{\pi'} \le
\pheight{\pi}$. We do this by induction on $n = \pheight{\pi}$.

If $n=0$, $\pi$ is just an axiom. But if $\Gamma \Sequent \Delta, !A,
!A$ is an axiom of~\Log{G3c}, so is $\Gamma \Sequent \Delta, !A$:
either some atomic~$!C$ is !!a{element} of both $\Gamma$ and~$\Delta$,
or $!A$ is atomic and !!a{element} of~$\Gamma$. The same reasoning
applies if the end-sequent is $!A, !A, \Gamma \Sequent
\Delta$.

If $n>0$, $\pi$ ends in some rule~$R$. If $!A$ is not principal in
this last inference, we can obtain !!a{proof}~$\pi'$ as in the proof
of \olref{lem:G3c-invert} by applying the inductive hypothesis to the
immediate sub-!!{proof}s of~$\pi$ and applying the same rule~$R$ to
the resulting !!{proof}(s). The inductive hypothesis is that if there
is !!a{proof} of !!{height}~$k<n$ of either $\Pi \Sequent \Lambda, !D,
!D$ or $!D, !D, \Pi \Sequent \Lambda$, then there is !!a{proof} of
!!{height}~$\le k$ of $\Pi \Sequent \Lambda, !D$ or $!D, \Pi \Sequent
\Lambda$, respectively. (Note that the inductive hypothesis applies if
the contexts or contracted !!{formula} are different from $\Gamma$,
$\Delta$, or~$!A$!)

For instance, suppose $R$ is $\RightR{\land}$
and $\pi$~ends in
\[
\AxiomC{}
\RightLabel{$\pi_1$}
\Deduce$\Gamma \fCenter \Delta', !B, !A, !A$
\AxiomC{}
\RightLabel{$\pi_2$}
\Deduce$\Gamma \fCenter \Delta', !C, !A, !A$
\RightLabel{\RightR{\land}}
\BinaryInf$\Gamma \fCenter \Delta', !B \land !C, !A, !A$
\DisplayProof
\]
where $\Delta = \Delta', !B \land !C$. The inductive hypothesis yields
!!{proof}s $\pi_1'$ and $\pi_2'$ of $\Gamma \fCenter \Delta', !B, !A$
and $\Gamma \fCenter \Delta', !C, !A$, respectively, with
$\pheight{\pi_1'} \le  \pheight{\pi_1}$ and $\pheight{\pi_2'} \le
\pheight{\pi_2}$. The !!{proof}~$\pi'$ is:
\[
\AxiomC{}
\RightLabel{$\pi_1'$}
\Deduce$\Gamma \fCenter \Delta', !B, !A$
\AxiomC{}
\RightLabel{$\pi_2'$}
\Deduce$\Gamma \fCenter \Delta', !C, !A$
\RightLabel{\RightR{\land}}
\BinaryInf$\Gamma \fCenter \Delta', !B \land !C, !A$
\DisplayProof
\]

If $!A$ is principal in the last inference, we apply the inversion
lemma (\olref{lem:G3c-invert}) to the immediate sub-!!{proof}s
of~$\pi$, then the inductive hypothesis, and then rule~$R$ again. For
instance, suppose $\pi$ proves $\Gamma \Sequent \Delta, !A, !A$, $!A
\ident (!B \land !C)$ and principal in the last inference. Then $\pi$
looks as follows:
\[
\AxiomC{}
\RightLabel{$\pi_1$}
\Deduce$\Gamma \fCenter \Delta, !B \land !C, !B$
\AxiomC{}
\RightLabel{$\pi_2$}
\Deduce$\Gamma \fCenter \Delta, !B \land !C, !C$
\RightLabel{\RightR{\land}}
\BinaryInf$\Gamma \fCenter \Delta, !B \land !C, !B \land !C$
\DisplayProof
\]
Apply \olref{lem:G3c-invert} to $\pi_1$ to get !!a{proof}~$\pi_1'$ of
$\Gamma \Sequent \Delta, !B, !B$. (It also yields !!a{proof} of
$\Gamma \Sequent \Delta, !C, !B$, but we don't need that.) Similarly,
by applying \olref{lem:G3c-invert} to~$\pi_2$ we obtain
!!a{proof}~$\pi_2'$ of $\Gamma \Sequent \Delta, !C, !C$. As the
inversion of \RightR{\land} is !!{height}-preserving,
$\pheight{\pi_1'} \le  \pheight{\pi_1} < \pheight{\pi}$ and
$\pheight{\pi_2'} \le  \pheight{\pi_2} < \pheight{\pi}$. Thus, the
inductive hypothesis applies to $\pi_1'$ and $\pi_2'$: there are
!!{proof}s~$\pi_1''$ and $\pi_2''$ of $\Gamma \Sequent \Delta, !B$ and
$\Gamma \Sequent \Delta, !C$, respectively, with $\pheight{\pi_1''}
\le \pheight{\pi_1'} \le \pheight{\pi_1}$ and $\pheight{\pi_2''} \le
\pheight{\pi_2'} \le \pheight{\pi_2}$. The !!{proof}~$\pi'$ then is:
\[
\AxiomC{}
\RightLabel{$\pi_1''$}
\Deduce$\Gamma \fCenter \Delta, !B$
\AxiomC{}
\RightLabel{$\pi_2''$}
\Deduce$\Gamma \fCenter \Delta, !C$
\RightLabel{\RightR{\land}}
\BinaryInf$\Gamma \fCenter \Delta, !B \land !C$
\DisplayProof
\]
with $\pheight{\pi'} \le \pheight{\pi}$.

We must take special care for the quantifier rules in the cases where
the !!{formula} is principal.

Suppose $!A \ident \lforall[x][!B(x)]$ is principal on the right. Then $\pi$ ends in
\[
\AxiomC{}
\RightLabel{$\pi_1$}
\Deduce$\Gamma \fCenter \Delta, \lforall[x][!B(x)], !B(a)$
\RightLabel{\RightR{\lexists}}
\UnaryInf$\Gamma \fCenter \Delta, \lforall[x][!B(x)], \lforall[x][!B(x)]$
\DisplayProof
\]
By \olref{lem:invert-quant}, there is !!a{proof}~$\pi_1'$ with
!!{height} $\le \pheight{\pi_1}$ of $\Gamma \fCenter \Delta, !B(c),
!B(a)$ for any~$c$ not occuring in $\Gamma \fCenter \Delta,
\lforall[x][!B(x)], !B(a)$. We may assume that $\pi_1'$ is regular; so
by \olref[qua]{lem:subst-regular} the !!{proof} $\Subst{\pi_1'}{a}{c}$
is !!a{proof} of $\Gamma \fCenter \Delta, !B(a),
!B(a)$ also of !!{height} $\le \pheight{\pi_1}$. Now the inductive
hypothesis applies and gives us !!a{proof}~$\pi_1''$ of $\Gamma
\fCenter \Delta, !B(a)$. Apply the \RightR{\lforall} rule to
get~$\pi'$:
\[
\AxiomC{}
\RightLabel{$\pi_1''$}
\Deduce$\Gamma \fCenter \Delta, !B(a)$
\RightLabel{\RightR{\lforall}}
\UnaryInf$\Gamma \fCenter \Delta, \lforall[x][!B(x)]$
\DisplayProof
\]
with $\pheight{\pi'} \le \pheight{\pi}$.

Now suppose $!A \ident \lexists[x][!B(x)]$
and principal on the right. Then $\pi$ looks as follows:
\[
\AxiomC{}
\RightLabel{$\pi_1$}
\Deduce$\Gamma \fCenter \Delta, \lexists[x][!B(x)], \lexists[x][!B(x)], !B(t)$
\RightLabel{\RightR{\lexists}}
\UnaryInf$\Gamma \fCenter \Delta, \lexists[x][!B(x)], \lexists[x][!B(x)]$
\DisplayProof
\]
The inductive hypothesis applies to~$\pi_1$ and there is
!!a{proof}s~$\pi_1'$ of $\Gamma \Sequent \Delta, \lexists[x][!B(x)],
!B$. (Note that no inversion lemma is needed here!) The
!!{proof}~$\pi'$ then is:
\[
\AxiomC{}
\RightLabel{$\pi_1'$}
\Deduce$\Gamma \fCenter \Delta,  \lexists[x][!B(x)], !B$
\RightLabel{\RightR{\lexists}}
\UnaryInf$\Gamma \fCenter \Delta, \lexists[x][!B(x)]$
\DisplayProof
\]
with $\pheight{\pi'} \le \pheight{\pi}$.
\end{proof}

\begin{prob}
Sketch the proof of \olref{prop:G3c-cont-adm} for \LeftR{\Contraction}.
Describe the cases where $!A \ident (!B \land !C)$ and $!A \ident (!B
\lif !C)$.
\end{prob}

\begin{prob}
Sketch the proof of \olref{prop:G3c-cont-adm} for \LeftR{\Contraction}
for the remaining quantifier cases, i.e., where $!A \ident
\lforall[x][!B(x)]$ and $!A \ident \lexists[x][!B(x)]$.
\end{prob}


\end{document}